\section{Related Work and Research Position}
\label{sec:related}
\subsection{Generative trajectories and reconstruction}
Diffusion models learn a denoising process, and implicit sampling exposes a deterministic trajectory suitable for inversion \cite{ho2020ddpm,song2021ddim}. Latent diffusion reduces the cost of high-resolution synthesis by operating in a learned image representation \cite{rombach2022ldm}. Flow matching and rectified flow learn continuous transport fields; transformer-based diffusion and scaled rectified-flow models provide architectures for joint image and text processing \cite{lipman2023flowmatching,liu2023rectifiedflow,peebles2023dit,esser2024sd3}. These foundations explain why a perturbation to a latent velocity or an attention output can influence several subsequent generation steps.

Real-image editing requires a useful source representation as well as a target instruction. SDEdit controls the amount of source destruction through noise injection \cite{meng2022sdedit}. Null-text inversion optimizes the unconditional embedding along a reconstruction trajectory, EDICT uses coupled transformations for invertibility, and PnP Inversion applies reconstruction offsets while separating source and target processes \cite{mokady2023nulltext,wallace2023edict,ju2024pnpinversion}. In rectified flows, RF-Inversion introduces endpoint-directed control, RF-Solver improves integration, and FireFlow reuses velocity information for efficient inversion \cite{rout2025rfinversion,wang2025rfsolver,deng2025fireflow}. TraceMemory uses a controlled source trajectory as a reference and measures the effect of adding a spatially allocated velocity residual during target generation.

\subsection{Direct flow editing and reference alignment}
FlowEdit constructs direct source-to-target transport from differences of source and target velocities evaluated along noised source and edited states \cite{kulikov2025flowedit}. CVC separates trajectory drift and semantic velocity differences; FlowAlign adds trajectory regularization; optimal-transport guidance supplies another route to source-aware flow editing \cite{li2025cvc,kim2025flowalign,lupascu2026otrf}. These methods make the evaluation state and reference construction central design choices.

DirectEdit saves adjacent inversion offsets, shifts the forward field-evaluation input, and combines source-value injection with mask-guided latent blending \cite{yang2026directedit}. SteerFlow combines reference-velocity alignment, trajectory interpolation, and adaptive masking derived from segmentation and source--target velocity differences \cite{dao2026steerflow}. Both are close comparisons for spatial source constraints. In the present method the target field remains global, while an additive correction is routed through a protection region and normalized explicitly. The experimental distinction concerns the location and magnitude of this residual and its interaction with a separate attention-memory channel.

\subsection{Attention and feature reuse}
Prompt-to-Prompt controls cross-attention to preserve spatial arrangements during prompt changes \cite{hertz2023prompttoprompt}. Plug-and-Play Diffusion Features transfers spatial features and self-attention to guide image translation \cite{tumanyan2023pnpfeatures}. MasaCtrl retrieves source keys and values using target queries, with mask-guided foreground/background retrieval for consistency \cite{cao2023masactrl}. These approaches establish attention reuse as a source-preservation mechanism.

Flow-based feature reuse differs in the features stored, the extraction state, and the receiving layers. RF-Edit shares source values; ReFlex extracts intermediate attention and residual features and adapts their injection; Stable Flow identifies influential layers for source-feature transfer; FluxSpace uses transformer features for semantic editing and localization \cite{wang2025rfsolver,kim2025reflex,avrahami2025stableflow,dalva2025fluxspace}. TraceMemory's local attention operation retains the target text contribution and the native image-attention mass while retrieving paired source image K/V. The comparison of all-query and ROI-query memory asks how coverage interacts with residual correction. This local algebraic preservation is defined at a fixed layer input; its final semantic effect is an empirical question.

\subsection{Spatial control and the meaning of a mask}
Blended Diffusion and Blended Latent Diffusion use spatial constraints to preserve image content during localized editing \cite{avrahami2022blended,avrahami2023blendedlatent}. DiffEdit infers an editing region from conditional model responses \cite{couairon2023diffedit}. GLIGEN introduces grounded spatial conditioning for generation, showing how explicit location information can control a generative model \cite{li2023gligen}. InstructPix2Pix learns instruction-conditioned editing from paired examples, while DiffusionCLIP and Imagic adapt models or representations to a requested change \cite{brooks2023instructpix2pix,kim2022diffusionclip,kawar2023imagic}. FreeDoM instead applies training-free energy guidance \cite{yu2023freedom}.

A protection mask and an edit mask serve different purposes. A shirt-print region identifies content to preserve; a studio-background edit region identifies content to change. Segment Anything and Grounding DINO provide segmentation and language-grounded detection components that can construct source-based regions \cite{kirillov2023sam,liu2024groundingdino}. Their outputs would be inputs to a future automatic protection pipeline. The experiments reported here use manual pattern regions. Their coordinates are also used by the local evaluation metric, and this annotation-assisted setting is retained explicitly throughout the results.

\subsection{Frequency, temporal allocation, and multi-turn editing}
FreqEdit injects high-frequency velocity information and compensates accumulated deviation for multi-turn editing \cite{liao2025freqedit}. Wavelet-guided semantic compensation extracts low-frequency same-state conditional differences as an auxiliary term for direct flow editing \cite{tang2026wavelet}. Multi-turn Consistent Image Editing combines control and attention adaptation across edits \cite{zhou2025multiturn}. These works motivate evaluating which quantity is filtered, where it is evaluated, and how its magnitude is controlled. Our spectral control filters the reference residual; its matched unfiltered counterpart isolates the effect of frequency shaping within the same spatial support.

\subsection{Evaluation and the specific research question}
CLIP provides image--text embeddings, and CLIPScore uses their similarity for reference-free image--text evaluation \cite{radford2021clip,hessel2021clipscore}. DINO and DINOv2 supply visual representations useful for content and structure comparisons \cite{caron2021dino,oquab2024dinov2}. LPIPS measures perceptual differences, and SSIM compares local structural information \cite{zhang2018lpips,wang2004ssim}. These measures address complementary outcomes. Exact logo characters, repeated objects, and a requested directional light require task-specific judgments in addition to embedding similarity.

Across this literature, source preservation is implemented through reconstruction, feature reuse, spatial constraints, or velocity guidance. The specific question studied here is how \emph{spatial support, applied correction magnitude, and memory-query coverage} interact within one fixed-reference editor. Matching the same coefficient does not match the injected RMS after masking, and changing attention coverage changes the trajectory on which the next residual is evaluated. The controlled studies quantify these interactions. Section~\ref{sec:extension} defines request-conditioned support. Appendix~\ref{sec:prospective-controls} separates area-aware and source-pool hypotheses from the completed experiments.

\begin{table*}[t]
\centering\small
\caption{Mechanism-level comparison of the nearest literature. This table records algorithmic relationships, not a measured performance ranking.}
\label{tab:nearest}
\begin{tabular}{p{.14\textwidth}p{.24\textwidth}p{.23\textwidth}p{.29\textwidth}}
\toprule
Method & Source constraint & Spatial / memory operation & Specific comparison for TraceMemory\\
\midrule
RF-Edit \cite{wang2025rfsolver} & Inversion trajectory and improved integration & Early source-value sharing & Source values versus paired K/V retrieval with retained target text contribution\\
MasaCtrl \cite{cao2023masactrl} & Source self-attention features & Source-pool masks and target-region combination & Full source pool with a query-output gate; correction controlled separately\\
ReFlex \cite{kim2025reflex} & Mid-inversion features & Attention and residual adaptation & Single-node cache versus seven reference captures and an explicit residual\\
DirectEdit \cite{yang2026directedit} & Adjacent inversion offsets & Input alignment, value injection, latent blending & Input alignment versus additive regional velocity correction\\
SteerFlow \cite{dao2026steerflow} & Reference alignment and interpolation & Adaptive edit masks and velocity steering & Edit-region velocity interpolation versus protection-region residual allocation\\
FlowEdit \cite{kulikov2025flowedit} & Randomly noised source states & Direct cross-state velocity differences & Inversion-free editing baseline under a shared input and evaluation protocol\\
TraceMemory & Controlled fixed trace & Regional residual plus separate paired K/V memory & Spatial support, magnitude normalization, and query coverage varied explicitly\\
\bottomrule
\end{tabular}
\end{table*}
